% Section 8.2, from lectures/L08/appendix/b_deferred_work.md.

\section{Deferred Work}
\label{c8:app:b}

A handler must be quick and must not sleep. Most real work is one or the other. This section is the
four places to put it, what each costs, and why one of them is the answer nearly every time now.

\subsection{The split, and why it has two names}
\label{c8:app:b1}

The kernel's traditional vocabulary is \textbf{top half} and \textbf{bottom half}. The top half is
the handler that runs in interrupt context: acknowledge the device, take the data out of it, and get
out. The bottom half is everything else, run later, in a context with fewer restrictions.

The names are old and slightly misleading, because ``bottom half'' also names one specific mechanism
(\code{local_bh_disable} still refers to it). Modern material says \textbf{hard IRQ handler} and
\textbf{deferred work}, and this section does too.

The question is not whether to split, it is \textbf{what to split at}. The dividing line is: the
hard handler does what only it can do, which is whatever must happen before the device is quiet and
whatever would be lost if it waited. Everything else is deferred.

For the lab's device that means: read \code{IRQ_STATUS}, acknowledge, drain the FIFO so no samples
are dropped, return. Formatting, copying to userspace, waking readers and anything that allocates
belongs on the other side.

\subsection{Four places to put it}
\label{c8:app:b2}

\begin{booktable}{@{}lll>{\hsize=0.58\hsize}L>{\hsize=1.42\hsize}L@{}}
  \thead{Mechanism} & \thead{Context} & \thead{May sleep} & \thead{Priority} & \thead{Use it} \\
  \midrule
  \textbf{Softirq} & Interrupt-ish & \textbf{No} & Fixed, very high & Almost never. Networking and timers own it \\
  \textbf{Tasklet} & Interrupt-ish & \textbf{No} & Fixed & \textbf{Deprecated.} Recognise it in old code \\
  \textbf{Workqueue} & Kernel thread & \textbf{Yes} & Schedulable & Work that sleeps and is not urgent \\
  \textbf{Threaded handler} & Kernel thread & \textbf{Yes} & Settable, per-IRQ & The default answer for a driver \\
\end{booktable}

\subsubsection{Softirqs}
\label{c8:app:b:softirqs}

A fixed, compile-time set of ten or so, defined in the kernel and not extensible by modules. They
run after the hard handler returns, still in a context that may not sleep, and they can run
concurrently with themselves on different CPUs. Networking and the timer wheel use them. \textbf{You
will not add one}, and the reason to know they exist is that they are what \code{local_bh_disable}
and \code{spin_lock_bh} are protecting against.

\subsubsection{Tasklets}
\label{c8:app:b:tasklets}

A softirq with a nicer interface, historically the standard bottom half for drivers. \textbf{They
are deprecated}, and the kernel's own documentation says so. Two reasons: they run at a fixed high
priority nobody can tune, and they still cannot sleep, so they solve only half the problem. Code
using them is not broken, it is old; new code should not add any.

\begin{ccode}
/* You will meet this. Do not write it. */
DECLARE_TASKLET(my_tasklet, my_func);
tasklet_schedule(&my_tasklet);
\end{ccode}

\subsubsection{Workqueues}
\label{c8:app:b:workqueues}

Work runs on a kernel thread, so it \textbf{may sleep}: allocate with \code{GFP_KERNEL}, take
mutexes, do I/O.

\begin{ccode}
static void my_work_fn(struct work_struct *work) { /* may sleep */ }
static DECLARE_WORK(my_work, my_work_fn);

/* from the hard handler */
schedule_work(&my_work);
\end{ccode}

\noindent \code{schedule_work} puts it on the shared system queue, which is convenient and means
your work waits behind everyone else's. A driver that cares creates its own with
\code{alloc_workqueue}, and one that cares about ordering uses a single-threaded one.

The catch is that \textbf{the same work item cannot be queued twice}. If it is already pending, a
second \code{schedule_work} does nothing, so a burst of interrupts produces one run. Sometimes that
is exactly right; when it is not, you need a queue of your own between the two halves, which is what
the FIFO in this course's driver is.

\subsection{What runs when}
\label{c8:app:b3}

Worth having straight, because the ordering explains a class of confusing latency.

\begin{console}
device asserts line
  -> CPU takes the interrupt, masks it
     -> hard IRQ handler runs          (atomic, fast, must not sleep)
  -> handler returns
     -> softirqs run                   (still atomic)
        -> tasklets are a softirq
     -> scheduler may now run something else
        -> workqueue thread runs when scheduled   (may sleep)
        -> threaded IRQ handler runs when scheduled (may sleep)
\end{console}

\noindent The hard handler is the only part guaranteed to run promptly. Everything below it runs
\emph{when the scheduler gets to it}, and how soon that is depends on load, priority and the
preemption model, which is Chapter~\ref{c12:ch}'s subject entirely.

\subsection{Threaded handlers}
\label{c8:app:b4}

The modern answer, and the one to reach for by default:

\begin{ccode}
request_threaded_irq(irq, qa_hard, qa_thread, IRQF_ONESHOT, "qa_irq", priv);
\end{ccode}

\noindent Two functions. \code{qa_hard} runs in interrupt context and returns
\code{IRQ_WAKE_THREAD}; \code{qa_thread} then runs in a dedicated kernel thread and may sleep.

You can see the thread:

\begin{console}
# ps | grep irq
   85 root      0:00 [irq/17-qa_irq]
\end{console}

\noindent It is an ordinary task, with a PID, that the scheduler can preempt, prioritise and account
for. That is the whole advantage over a workqueue: \textbf{the deferred work has an identity}. You
can give it a real-time priority, pin it to a CPU, and see it in \code{top} when it misbehaves.

\textbf{\code{IRQF_ONESHOT} is not optional for a level-triggered device.} It keeps the interrupt
masked from the moment the hard handler returns until the thread has finished. Without it the line
is unmasked as soon as the hard handler returns, and since a level-triggered device is still
asserting until acknowledged, the hard handler is re-entered immediately, which is
\secref{c8:app:a6}. Either the hard handler acknowledges before returning, or you use
\code{IRQF_ONESHOT}.

Passing \code{NULL} as the hard handler is allowed and means ``wake the thread for every
interrupt''; \code{IRQF_ONESHOT} is then mandatory.

\subsubsection{Measured, on this target}
\label{c8:app:b:measured-on-this-target}

The lab's module both ways, at a one-millisecond period, over two seconds:

\begin{narrowtable}{@{}lrrr@{}}
   & \thead{Elapsed} & \thead{Interrupts serviced} & \thead{Shortfall against ideal} \\
  \midrule
  Hard handler only & 2.11 s & 1789 & 15.2\% \\
  Threaded handler & 2.10 s & 1792 & 14.7\% \\
\end{narrowtable}

\noindent \textbf{The threaded version is not slower}, which surprises people who expect the extra
context switch to cost. For work of this size the scheduler absorbs it, and what you bought is a
handler that may sleep and a thread you can prioritise. Under \code{PREEMPT_RT} almost every handler
becomes one of these whether you asked or not, which is why Chapter~\ref{c12:ch} revisits this
table.

\subsection{Choosing}
\label{c8:app:b5}

In order of preference:
\begin{enumerate}
  \item \textbf{Do it in the hard handler}, if it is a few register accesses and cannot sleep.
  \item \textbf{Threaded handler}, if it needs to sleep or is long enough to matter. This is the
    default.
  \item \textbf{Workqueue}, if the work is not per-interrupt: something periodic, or triggered from
    several places, or genuinely low priority.
  \item \textbf{Tasklet}, never in new code.
  \item \textbf{Softirq}, not available to you.
\end{enumerate}
Two questions decide between 2 and 3. \emph{Is this work one-per-interrupt?} If yes, a threaded
handler already has the plumbing. \emph{Does it need its own priority?} If yes, a threaded handler
has one and the shared workqueue does not.

\subsection{What deferring does not fix}
\label{c8:app:b6}

It moves work; it does not make it free, and three things stay broken.

\textbf{The device still has to be serviced promptly.} \code{qa-dev} has a FIFO of sixteen samples
and sets \code{OVERRUN} when it overflows. Deferring the drain does not slow the device down; it just
means the FIFO fills while you wait. If the hard handler does not empty it, no amount of deferred
work will recover the samples that were dropped.

\textbf{A high interrupt rate is still a high interrupt rate.} At a hundred-microsecond period the
lab's module services about 13,000 interrupts in two seconds against an ideal of 21,300, a
\textbf{39\% shortfall}, and no arrangement of deferred work improves that: the cost is in taking
the interrupt at all.

What \emph{does} happen at that rate is visible in the driver's own two counters:

\begin{console}
period=1000000   hard=1813   samples=1813     <- one sample per interrupt
period=100000    hard=13708  samples=14389    <- more samples than interrupts
\end{console}

\noindent At 1 ms the handler keeps up exactly. At 0.1 ms it does not, and the device's
level-triggered line does the only sensible thing: a sample arriving while the line is still
asserted does not generate a second interrupt, it simply joins the FIFO that the handler is about to
drain. \textbf{The interrupts coalesced, and no data was lost.}

That is level triggering paying for itself, and it is the answer to why a device under load
generates fewer interrupts than its nominal rate rather than more. Push harder and the FIFO fills,
\code{OVERRUN} sets, and samples do start being lost; at that point the answer is hardware, in the
shape of a bigger FIFO, explicit coalescing, or DMA.

\textbf{The total work is unchanged.} Deferring improves \emph{latency for everything else} and
worsens \emph{latency for this device's data}, which is a trade rather than a win. Knowing which of
the two you are optimising is the point, and Chapter~\ref{c12:ch} is where it gets measured rather
than argued about.
